\section{List of Textbooks}
\label{sec-2}

\section*{Introduction}
This is a list of textbooks suitable for fourth-year undergraduate seminars and master's seminars in Japan. 
I used AI transfer (chatGPT), so some parts of the translation may be awkward, although it is much better than my own translation.
%Most of the books listed here are in Japanese, but I also include English editions when available.

\subsection{Books for undergraduate and master's seminars}
This list is written with fourth-year undergraduate seminars and master's seminars in mind.
%ただせっかくなので, 学部2$\sim$3年生向けの本もリストアップした. 
The choice of a textbook for a seminar depends on the student's situation. I roughly divide the cases into the following three types:
\begin{enumerate}[label=$(\arabic*)$]
  \setlength{\parskip}{0cm} % 段落間
  \setlength{\itemsep}{0cm} % 項目間
 \item Students in a fourth-year undergraduate seminar who do not plan to enter a master's program.
\item Students in a fourth-year undergraduate seminar who do plan to enter a master's program.
\item First-year master's students in a master's seminar.
\end{enumerate}

(1). If you do not plan to enter a master's program in mathematics, then you do not need to worry too much about which book to read in a fourth-year undergraduate seminar. You may read one of the books listed in (2) below, or, if those books are too difficult, you may choose an easier one.  For examples of such books, see Section \ref{subsec-seminar-23}.

(2). In a fourth-year undergraduate seminar, I recommend the following books to students who plan to enter a master's program if you plan to study Algebraic geometry, Complex geometry, or Several complex variables.

\begin{tcolorbox}[mybox]
\begin{itemize}[left=0pt]
  \setlength{\parskip}{0cm} % 段落間
  \setlength{\itemsep}{8pt} % 項目間
\item If you want to specialize in algebraic geometry, read Hartshorne's textbook. 
The reason is that this book allows you to learn schemes, sheaves, and sheaf cohomology.
\item If you want to specialize in complex geometry or several complex variables, read textbook by Wells, Huybrechts, or Horikawa. 
The reason is that these books allow you to learn sheaves, cohomology, Hodge theory (harmonic integrals), the Kodaira vanishing theorem, and the Kodaira embedding theorem. (Horikawa's textbook  is an exception in this respect.)
\item If you want to study several complex variables, read textbook of Nakano or Forstneric. 
About half of Nakano's book concerns complex geometry, and the other half concerns $L^2$ theory. Forstneric's book treats Oka manifolds. I expect that Oka manifolds may become established as a new area of several complex variables.
\end{itemize}
\end{tcolorbox}
%Of course, other books are also possible, but the textbooks listed above are safe choices (although it is not clear whether Forstneric is a safe choice). Here, by ``safe'' I mean that, even if you later enter a master's program under a professor other than me, you should be able to continue without serious difficulty.

(3). In a master's seminar, you will study more specialized material. The following three areas are the ones that I can cover:
\begin{enumerate}
  \setlength{\parskip}{0cm} % 段落間
  \setlength{\itemsep}{0pt} % 項目間
  \item Complex algebraic geometry arising from algebraic geometry and complex analytic spaces.
\item Complex differential geometry, such as the existence of K\"ahler--Einstein metrics, which is a classical and important subject at Osaka University.
\item Complex analytic geometry arising from the theory of domains and $L^2$ methods in several complex variables.
\end{enumerate}
My own specialty is (1). At present, however, this is a field in which there seem to be relatively few problems that master's or doctoral students can realistically solve. 
Since (2) and (3) are not my specialties, I cannot teach all details. However, if you intend to enter a doctoral program, you will need to go beyond your advisor. Therefore, I would be happy if you study (2) or (3) under my supervision.

For a master's seminar, I recommend the following books.
\begin{tcolorbox}[mybox]
\begin{itemize}[left=0pt]
  \setlength{\parskip}{0cm} % 段落間
  \setlength{\itemsep}{8pt} % 項目間
\item If you want to study algebraic geometry in (1) or complex analysis in (3), read Demailly's book \emph{Analytic Methods in Algebraic Geometry}. It takes about one year to read this book. (If you want to study complex analysis, it may be better to choose a slightly easier book, such as Berndtsson's book.) After reading this book, you will read the paper of Hacon--Popa--Schnell.  
\item If you want to study differential geometry in (2), read Szekelyhidi's book. You may read the whole book, or you may stop the later chapters somewhat early and then move on to papers. 
\item If you want to study several complex variables outside (1)--(3), then study Oka manifolds. In that case, you will read Forstneric's textbook or related papers.
\end{itemize}
\end{tcolorbox}

\vspace{8pt}
Interesting books and good textbooks are not necessarily the same. Hartshorne's book  is good, but whether it is interesting is another matter. 
Grauert--Remmert's book  is extremely interesting to read now, but I do not recommend it as a textbook. (I will explain the reason later.) I also use the expressions ``read as the main text'' and ``read as a supplementary text'' in the following sense.
\begin{itemize}
  \setlength{\parskip}{0cm} % 段落間
  \setlength{\itemsep}{0pt} % 項目間
\item ``Read as the main text'' means that the book is intended to be read in a seminar. Such books usually cover the knowledge needed for research, for example in a master's program.
\item ``Read as a supplementary text'' means that the book is suitable for reading alone or in an informal seminar with classmates. Such books treat more specialized material or material useful for research, but they are less suitable than the main texts as the central book of a seminar.
\end{itemize}
In one year, the number of books that one can read seriously is probably one main text and zero to two supplementary texts, that is, about one to three books in total. 
Therefore, the books that you read seriously should preferably be safe choices. On the other hand, more specialized or unusual topics may turn out to be useful later, and this is partly a matter of luck.

Finally, I quote the following words from Ihara's book.
\begin{tcolorbox}[mybox]
Do not be misled by senior students or peers who say things such as, ``I read one book every month.'' People who cannot grasp what is essential tend to make numerical evaluations too quickly. I think that becoming a mathematician also means aiming to become a person who can understand deeper values.\\
(Original Japanese text) 「自分は一ヶ月に一冊読んでいる」などという先輩や仲間に惑わされないように。
本質的なところを感じとれない人の方が、すぐ数値的な評価をしたがる。数学者に
なる、ということは、より深い価値がわかる人をめざす、ということでもあると思
います。
\end{tcolorbox}

\subsection{Difficulty levels of the books}
The star marks indicate the difficulty level. The levels are as follows.\footnote{The names of the difficulty levels are inspired by the Touhou Project and the Fire Emblem series.} 
\vspace{-8pt}
\begin{itemize}
  \setlength{\parskip}{0cm} % 段落間
  \setlength{\itemsep}{4pt} % 項目間
\item[★] Easy. For second- or third-year undergraduates. In a fourth-year undergraduate seminar, only students who do not plan to enter a master's program may choose books at this level.
\item[★★] Normal. For fourth-year undergraduates. Students who plan to enter a master's program should choose books at this level or above.
\item[★★★] Hard. For first-year master's students. Roughly speaking, if you can read a book at this level in your first year of a master's program, then you should be able to start reading papers afterwards.
\item[★★★★] Lunatic. For second-year master's students or above. Some textbooks and papers at this level are beyond even Iwai's ability to read.
\end{itemize}

The comments are based on my memories of reading these books when I was a student. The contents of the books may have changed in later editions. My memory may also be inaccurate, and many of the comments are personal opinions. Please keep this in mind.

In preparing this list, I referred to the following sources. Most of these references are written in Japanese.
\vspace{-8pt}
\begin{itemize}[left=0pt]
  \setlength{\parskip}{0cm} % 段落間
  \setlength{\itemsep}{4pt} % 項目間
\item Osamu Fujino. Literature guide  \url{https://www.math.kyoto-u.ac.jp/~fujino/ag-ref.pdf}
\item Toshiyuki Katsura, Yoshito Ohta, and Osami Yamamoto. 
  A practical guide to the literature on algebraic geometry (standard textbooks and classical books) \url{https://www.jstage.jst.go.jp/article/bjsiam/14/3/14_KJ00003509979/_article/-char/ja/}
\item Shinichiro Matsuo. Recommended books   \\ \url{https://www.math.nagoya-u.ac.jp/~shinichiroh/2018/01/20/recommended-texts.html}
\item Nobuhiro Honda. Textbooks and reference books for seminars of fourth-year undergraduates and graduate students  \\ \url{https://www.honda.math.sci.titech.ac.jp}
\item Department of Mathematics, The University of Tokyo. PDF introducing the research areas of faculty members \\ \url{https://www.ms.u-tokyo.ac.jp/kyoumu/kenkyubunya.html}
\end{itemize}

I mention one useful information. Springer Link \url{https://link.springer.com} allows you to download PDFs of Springer books for free, depending on the university to which you belong.

\subsection{Books that everyone should read}

\textbf{Yasutaka Ihara. "志学数学 -研究の諸段階 発表の工夫" Maruzen Publishing (in Japanese)} \vspace{-6pt}

Difficulty: none \vspace{-6pt}

Comment: This is a book that everyone should read, regardless of their field. If you have not read it, I recommend reading it immediately. Professor Kawahigashi also wrote, ``Go buy this book and read it right now.'' in his review. I completely agree.

Professor Kawahigashi's book \emph{数学者の思案} by Iwanami Shoten (in Japanese) is also interesting. I bought the Kindle edition. His website is also interesting. I recommend reading the following pages in Japanese:
\begin{itemize}[left=0pt]
  \setlength{\parskip}{0cm} % 段落間
  \setlength{\itemsep}{0cm} % 項目間
\item セミナーの準備のしかたについて \url{https://www.ms.u-tokyo.ac.jp/~yasuyuki/sem.htm}
\item どうでもよい記事 \url{https://www.ms.u-tokyo.ac.jp/~yasuyuki/misc.htm}
\end{itemize}
どうでもよい記事 are something I sometimes want to reread. They also contain much useful information. It is almost surprising that they are available for free.

\subsection{Textbooks in algebraic geometry}

If you study algebraic geometry, then through a fourth-year undergraduate seminar you need to understand schemes, sheaves, and sheaf cohomology. To understand schemes, you also need to know some commutative algebra. Therefore, by the time the fourth-year seminar begins, you should have studied at least a certain amount of commutative algebra, roughly the material in Atiyah--Macdonald.

\subsubsection{For fourth-year undergraduate seminars}

\textbf{Robin Hartshorne. "Algebraic Geometry" Graduate Texts in Mathematics (GTM, volume 52), Springer} \vspace{-6pt}

Difficulty: ★★$\sim$★★★  \vspace{-6pt}

Comment: There is only one reason why I recommend this book: many papers cite it. Of course, there are algebraic geometry textbooks that are easier to read than this one. However, if you have no particular preference, I think it is better to read this book.

If you read it in one year, read Chapters 2 and 3. The reputation of Chapter 2 on schemes is not very good, but this cannot be helped. Chapter 3 on cohomology is the most interesting. Chapter 1 on classical algebraic geometry can be read beforehand, or it can be read later if necessary. Chapter 4 on curves is useful, so if you finish Chapter 3 and still have time, you may read it. For Chapter 5 on surfaces, it is not necessary to use this book, although reading it is also fine.
\vspace{8pt}

\textbf{Masayoshi Miyanishi. "代数幾何学" Shokabo (in Japanese)}  \vspace{-6pt}

Difficulty: ★★$\sim$★★★  \vspace{-6pt}
 
Comment: Chapter 1 gives a fairly detailed account of sheaf theory, which is a good point, although it is somewhat hard to read. The part on schemes is standard. The final part on surface theory seems to reflect Professor Miyanishi's own field. 
\vspace{8pt}

\textbf{Kenji Ueno. "代数幾何" Iwanami Shoten (in Japanese)}\vspace{-6pt}

Difficulty: ★★$\sim$★★★  \vspace{-6pt}

Comment: I have not read this book. Judging from the table of contents, it seems to be a careful textbook.
\vspace{8pt}

\textbf{David Mumford. "Algebraic Geometry. I. Complex projective varieties" Springer.}\vspace{-6pt}

Difficulty: ★★$\sim$★★★  \vspace{-6pt}

Comment: This book is recommended by Professor Fujino and Professors Katsura and others. I have not read it. According to Professor Fujino, ``This book contains neither scheme theory, complex manifold theory, nor sheaf cohomology, but it is an excellent book that allows one to encounter the essence of algebraic geometry at an early stage.''

I think that, if one studies algebraic geometry, one should become familiar with schemes, sheaves, and cohomology in the fourth undergraduate year. For this reason, I do not recommend this book for a fourth-year undergraduate seminar. 
There seems to be a second volume, which treats schemes.
\vspace{8pt}

\textbf{David Mumford. The Red Book of Varieties and Schemes}\vspace{-6pt}

Difficulty: ★ $\sim$ ★★ \vspace{-6pt}

Comment: This book is commonly called the ``Red Book.'' I read about half of it in an informal seminar when I was a third-year undergraduate. I think it is a good book for getting a feeling for the subject. For the same reason as above, I do not recommend it for a fourth-year undergraduate seminar.
\vspace{8pt}

\textbf{Qing Liu. "Algebraic Geometry and Arithmetic Curves" Oxford Graduate Texts in Mathematics. Oxford University Press}

Difficulty: ★★$\sim$★★★  \vspace{-6pt}

Comment: My impression is that this book is read mainly by people in number theory.
\vspace{8pt}

\textbf{Ulrich G\"ortz , Torsten Wedhorn. "Algebraic Geometry I: Schemes" Springer}

Difficulty: ★★$\sim$★★★  \vspace{-6pt}

Comment: This book is long. If you like category theory, it may be a good book from which to learn algebraic geometry. The second volume was published recently.
\vspace{8pt}

\textbf{Shigeru Iitaka. "代数幾何学" Iwanami Lectures on Fundamental Mathematics. Iwanami Shoten  (in Japanese)} \vspace{-6pt}

\textbf{Sigeru Iitaka. "Algebraic Geometry
An Introduction to Birational Geometry of Algebraic Varieties"  (GTM, volume 76) Springer} \vspace{-6pt}

Difficulty: ★★$\sim$★★★  \vspace{-6pt} 

Comment: This is the book written by Professor Iitaka for the Iwanami Lectures on Fundamental Mathematics. When I was a fourth-year undergraduate, I read the English edition under Professor Shunsuke Takagi. The Japanese edition is too difficult for me to read, so I recommend reading the English edition.

As the title suggests, much of the content is connected to birational geometry. The first half contains basic material, but the second half contains many rather specialized topics, such as Grothendieck's d\'evissage theorem, Fujita's $\delta$-genus, and the Iitaka conjecture. It is probably better to glance through this book after studying algebraic geometry elsewhere than to learn the subject from this book. I also used this book in my research in 2021.
\vspace{8pt}

\textbf{"Elements de geometrie algebrique"} \vspace{-6pt}

\textbf{"Seminaire de Geometrie Algebrique du Bois Marie"} \vspace{-6pt}

Difficulty: ★★$\sim$★★★★  \vspace{-6pt} 

Comment: These are EGA and SGA. EGA has about 1500 pages, and SGA has about 6500 pages. They are too long, and moreover they are in French. Since I was an undergraduate, I have wondered, ``Who can read these?'' Yet even now it seems that many people read them. My impression is that number theorists read them. Algebraic geometers seem to read Hartshorne's textbook.

The famous Grothendieck--Riemann--Roch theorem appears in SGA6. I have used GRR, but I have never seen its proof. Also, when I recently looked up Grothendieck universes, I found that they also appear in SGA. Since even very set-theoretic material appears there, SGA seems to contain everything.

In fact, I wonder whether algebraic geometers before Hartshorne's book appeared read EGA. Professor Ohsawa once said, ``I could not understand EGA, but I could read Professor Nakano's book on algebraic geometry (Shigeo Nakano, \emph{Introduction to Algebraic Geometry}, (in Japanese)).''

\vspace{8pt}

\textbf{Shigeru Mukai. "モジュライ理論" Iwanami Shoten (in Japanese)}
\vspace{-6pt}

\textbf{Shigeru Mukai.  "An Introduction to Invariants and Moduli" Cambridge University Press}

Difficulty: ★★$\sim$★★★  \vspace{-6pt}

Comment: I am curious about this book, so I would like someone to read it.

\subsubsection{For master's seminars}

\textbf{Janos Kollar, Shigefumi Mori. "双有理幾何学" Iwanami Shoten} \vspace{-6pt}

\textbf{Janos Kollar, Shigefumi Mori. "Birational Geometry of Algebraic Varieties" Cambridge Tracts in Mathematics, Series Number 134. Cambridge University Press} \vspace{-6pt}

Difficulty: ★★★$\sim$★★★★\vspace{-6pt} 

Comment: If you study birational geometry, you will study this book in the first year of your master's program. This is because many papers cite Hartshorne and Koll\'ar--Mori by saying, for example, ``We employ the standard notation and conventions in [Har77, KM98].'' I have written this sentence myself.

In this field, one cannot avoid birational geometry and the minimal model program (MMP), so I recommend reading this book, even if only roughly. I have read it only roughly, and for that reason I still have to open this book from time to time.
 \vspace{8pt}
 
 \textbf{Yujiro Kawamata, Katsumi Matsuda, Kenji Matsuki. "Introduction to the Minimal Model Problem" Adv. Stud. Pure Math., 1987. } \vspace{-6pt}

Difficulty: ★★★$\sim$★★★★\vspace{-6pt} 

Available online:
\url{https://projecteuclid.org/ebooks/advanced-studies-in-pure-mathematics/Algebraic-Geometry-Sendai-1985/chapter/Introduction-to-the-Minimal-Model-Problem/10.2969/aspm/01010283} 
\vspace{-6pt}

Comment: Personally, I like this better than Koll\'ar--Mori. I like the programming-like diagrams that were probably written by Professor Matsuki. It is commonly called KMM. Incidentally, in algebraic geometry and birational geometry, there are at least three papers called ``KMM.''; Kawamata-Matsuda-Matsuki, Kollar-Miyaoka-Mori, and Keel-Matsuki-Mckernan.
\vspace{8pt}

\textbf{Robert Lazarsfeld. "Positivity in Algebraic Geometry I, II" Springer} \vspace{-6pt}

Difficulty: ★★★$\sim$★★★★\vspace{-6pt} 

Comment: If you want to understand Demailly's book \emph{Analytic Methods} algebraically, read this book. My impression is that this book is read by people in positive characteristic, especially students in Professor Shunsuke Takagi's research group. To understand multiplier ideal sheaves algebraically, you will read Part III of this book as a supplementary text. I also benefited from the discussion of vector bundles in Part II in my research.

See also Professor Shunsuke Takagi's review in Japanese. (\url{https://www.ms.u-tokyo.ac.jp/~stakagi/academic/review_Lazarsfeld.pdf})
\vspace{8pt}

\textbf{Arnaud Beauville. "Complex Algebraic Surfaces" London Mathematical Society Student Texts, Series Number 34. Cambridge University Press} \vspace{-6pt}

Difficulty: ★★★$\sim$★★★★\vspace{-6pt} 

Comment: If you want to learn surface theory quickly, this is the book.
\vspace{8pt}

\textbf{Daniel Huybrechts, Manfred Lehn. "The Geometry of Moduli Spaces of Sheaves" Cambridge University Press} \vspace{-6pt}

Difficulty: ★★★$\sim$★★★★\vspace{-6pt} 

Comment: If you want to learn moduli of semistable sheaves and related topics, read this book.
 \vspace{8pt}
 
 \textbf{Caucher Birkar, Paolo Cascini, Christopher D. Hacon, James McKernan. "Existence of minimal models for varieties of log general type" Journal of American Mathematical Society, 2010.} \vspace{-6pt}

Difficulty: ★★★★\vspace{-6pt} 

Comment: This paper is commonly called BCHM. Usually, mathematical papers do not appear in search results by their common abbreviations, but this paper is so famous that if you search for ``BCHM,'' it appears at the top. I am curious about it and would like someone to read it, since I have only cited its results.
\vspace{8pt}

\textbf{Yujiro Kawamata. "高次元代数多様体論" Iwanami Shoten (in Japanese)} \vspace{-6pt}

\textbf{Yujiro Kawamata. "Algebraic Varieties: Minimal Models and Finite Generation" Cambridge University Press} \vspace{-6pt}

Difficulty: ★★★★\vspace{-6pt} 

Comment: This is Professor Kawamata's expository book on BCHM. I bought it, but I have not read it. The English edition was translated by Professor Chen Jiang.
\vspace{8pt}

\textbf{Osamu Fujino. "Foundations of the Minimal Model Program" Mathematical Society of Japan Memoirs. Mathematical Society of Japan} \vspace{-6pt}

Difficulty: ★★★★\vspace{-6pt} 

Comment: We in geometry and analysis can only treat KLT (Kawamata log terminal) singularities, because KLT means $L^2$ integrability. Algebraic geometers, however, can treat LC (log canonical) singularities. This book proves vanishing theorems for LC objects and even more general quasi-log schemes, and it also proves the basic theorems of the MMP. Therefore, if you want to understand vanishing theorems algebraically, you will read this book. I benefited from this book when citing Professor Fujino's papers in my research.

\vspace{8pt}

\textbf{Osamu Fujino. "Iitaka Conjecture An Introduction" Springer} \vspace{-6pt}

Difficulty: ★★★★\vspace{-6pt} 

Comment: This book covers the results of Viehweg, Fujino, and others on the Iitaka conjecture. If you want to understand algebraically the positivity of direct image sheaves, such as the work of P\u{a}un--Takayama and Hacon--Popa--Schnell, this is the book. I benefited greatly from it in my first research project. Viehweg's covering trick appears very often.

See also Professor Fujino's 2022 Nagoya University lecture notes in Japanese \url{https://www.math.kyoto-u.ac.jp/~fujino/sonota.html}.
\vspace{8pt}

\textbf{Kenji Ueno. "Classification Theory of Algebraic Varieties and Compact Complex Spaces" Springer}  　\vspace{-6pt} 

Difficulty: ★★ $\sim$ ★★★★  　\vspace{-6pt} 

Comment: This is both a book on algebraic geometry and a book on several complex variables. It probably treats classification theory before the framework of the MMP was established. (For details, see also Professor Fujino's 2022 Nagoya University lecture notes in Japanese \url{https://www.math.kyoto-u.ac.jp/~fujino/sonota.html}.) Papers on the Iitaka conjecture cite this textbook very often. I introduce it because I often refer to it. If you are interested, you may read it casually as a supplementary text.
\vspace{8pt}

\textbf{Yoichi Miyaoka , Thomas Peternell. "Geometry of Higher Dimensional Algebraic Varieties" Springer} \vspace{-6pt}

Difficulty: ★★★★\vspace{-6pt} 

Comment: This is partly an advertisement. The discussion of Chern classes related to Professor Miyaoka's research is interesting, but not many people seem to know this book. It is a good book. I benefited from it in my research on the second Chern class.

\textbf{Christian Schnell. "A graduate course on the Generic vanishing theorem" }  \vspace{-6pt} 

Difficulty: ★★★★\vspace{-6pt} 

Available online: \url{https://www.math.stonybrook.edu/~cschnell/pdf/notes/generic-vanishing.pdf} \vspace{-6pt} 

Comment: These are lecture notes by Professor Schnell. They treat concrete applications of derived categories and the Fourier--Mukai transform to birational geometry. These are quite powerful tools, so I would like you to read these notes.
\vspace{8pt}

\textbf{Janos Kollar. "Rational Curves on Algebraic Varieties" Springer} \vspace{-6pt}

\textbf{Olivier Debarre. "Higher-Dimensional Algebraic Geometry" Springer} \vspace{-6pt}

Difficulty: ★★★$\sim$★★★★\vspace{-6pt} 

Comment: If you want to become a user of rational curves, for example in the study of Fano varieties, you will read these books. If you want to read these books in a master's seminar, you should read them not with me but with Professor Fujita. Debarre is easier to read, while Koll\'ar is harder. However, Koll\'ar carefully constructs Hilbert schemes and Chow schemes. I also feel that Koll\'ar is cited more often.

I have only looked at the basics of RC (rationally connected) varieties. Since I have not mastered the material at all, every time I read a paper I wonder, ``What is $\mathrm{RatCurve}^{n}$?'' and open Koll\'ar's book. Arguments specific to rational curves appear very often in papers on Fano varieties and RC varieties. For this reason, Professor Hiromichi Takagi's survey in Japanese \url{https://repository.kulib.kyoto-u.ac.jp/items/4e7bdf98-f0fe-4c8b-bb1b-2602841ea5af} is truly helpful.
 \vspace{8pt}
 
 \textbf{Chenyang Xu. "K-stability of Fano Varieties" Cambridge University Press} \vspace{-6pt}

Difficulty: ★★★★\vspace{-6pt} 

Comment: It is a book that summarizes algebraic aspects of $K$-stability.
 \vspace{8pt}
 

\subsection{Textbooks in complex geometry}
Complex geometry covers a wide range of topics, so, unlike algebraic geometry, the contents that one must understand vary depending on the field.

Personally, I think it is enough if one can understand sheaves, cohomology, Hodge theory (harmonic integrals), the Kodaira vanishing theorem, and the Kodaira embedding theorem in a fourth-year undergraduate seminar. Probably the best method is to choose one textbook, understand its contents, and then gradually expand the range of material that one needs to understand.

My personal view is that it is relatively easy to learn complex geometry after learning algebraic geometry, but it is difficult to learn algebraic geometry after learning complex geometry. 
This is because abstract theories such as scheme theory in algebraic geometry are hard to understand unless one spends enough time reading them carefully.\footnote{Even without scheme theory, one can instead learn complex analytic spaces, but learning complex analytic spaces is harder than learning scheme theory, because the arguments are quite complicated.} 
Therefore, if one learns algebraic geometry after complex geometry, one has to learn it in a more practical way. More concretely, one learns by using tools from algebraic geometry and birational geometry. For example, one translates algebraic-geometric terminology such as ample and psef into complex-geometric terminology such as positive and having a semipositive curvature current.

If you specialize in complex geometry, you cannot avoid algebraic geometry and birational geometry. It is quite common for the minimal model program to appear suddenly when reading papers. This cannot be helped.

\subsubsection{For fourth-year undergraduate seminars}

\textbf{Raymond O. Wells. "Differential Analysis on Complex Manifolds" Springer}  \vspace{-6pt}

Difficulty: ★★　\vspace{-6pt}

Comment: When I was an undergraduate, this book was the standard textbook for complex geometry. I think it is a good book that covers the material one should learn.

However, Chapter 4 on pseudo-differential operators is not very well regarded. One should either completely ignore that part or learn it from another textbook. (For other textbooks, see Professor Shinichiro Matsuo's page in Japanese.) If you choose the former, then you probably like algebraic-style complex geometry; if you choose the latter, then you probably like analytic-style complex geometry. I belonged to the former type, but either choice is fine.
\vspace{8pt}

\textbf{Daniel Huybrechts. "Complex Geometry: An Introduction" Universitext, Springer} \vspace{-6pt}
 
Difficulty: ★★ \vspace{-6pt}

Comment: I have not read this book. However, after flipping through it, I thought that it looked quite good. What surprised me was that sheaf theory is placed in the appendix. Since sheaves are indispensable in complex geometry, one should study the appendix as well.
\vspace{8pt}

\textbf{Eiji Horikawa. "複素代数幾何学入門" Iwanami Shoten (in Japanese)}

Difficulty: ★★$\sim$★★★ \vspace{-6pt} 

Comment: This book treats algebraic geometry over the complex numbers from a complex-analytic viewpoint. I highly recommend it. If you want to study complex algebraic geometry while completely ignoring scheme theory, it may be good to study from this book. It is possible to read this book and then proceed toward complex algebraic geometry.

Up to the part on Riemann surfaces, except for the part on divisors, the contents are similar to those of traditional books.
 I think this is almost the only Japanese book that treats the final part on degenerations of elliptic curves.
  If you can read it to the end, you will acquire enough algebraic geometry to compete with algebraic geometers. The unfortunate point is that it does not contain Hodge theory (harmonic integrals), the Kodaira vanishing theorem, or the embedding theorem. You need to supplement these parts with Wells or Huybrechts.
\vspace{8pt}

\textbf{Shoshichi Kobayashi. "複素幾何" Iwanami Shoten  (in Japanese)}  \vspace{-6pt} 

Difficulty: ★★$\sim$★★★ \vspace{-6pt} 

Comment: If you want to learn complex geometry in Japanese, this may be the book. The part on the second fundamental form of vector bundles was extremely useful in my research. Of course, it is fine to study from this book. However, if you read a book in a seminar, I think it is better to read a book in English.
\vspace{8pt}

\textbf{Christian Schnell. "A graduate course on Complex manifolds" }  \vspace{-6pt} 

Difficulty: ★★$\sim$★★★ \vspace{-6pt} 

Available online: \url{https://www.math.stonybrook.edu/~cschnell/pdf/notes/complex-manifolds.pdf} \vspace{-6pt} 

Comment: These are lecture notes by Professor Schnell. I only glanced at them, but they also look good. The material one wants to learn is summarized concisely in 136 pages, and there are many algebraic-style examples. My impression is that Professor Schnell's papers and books are often easy to understand. 
\vspace{8pt}

\textbf{Claire Voisin. "Hodge Theory and Complex Algebraic Geometry I, II" Cambridge University Press}  \vspace{-6pt} 

Difficulty: ★★$\sim$★★★ \vspace{-6pt} 

Comment: The beginning is similar to traditional textbooks, but later the book moves on to topics such as VHS (variation of Hodge structure). I think this is one of the few books from which one can learn these topics from a geometric viewpoint. In this respect, the book is worth reading. 
\vspace{8pt}

\textbf{Kunihiko Kodaira. "複素多様体論" Iwanami Shoten  (in Japanese)}\vspace{-6pt}

\textbf{Kunihiko Kodaira "Complex Manifolds and Deformation of Complex Structures" Springer} \vspace{-6pt}

Difficulty: ★★$\sim$★★★ \vspace{-6pt} 

Comment: The first half is similar to traditional books, while the second half treats Kodaira--Spencer deformation theory. The English edition gives a more detailed proof of the deformation theory. The material in the second half is interesting, but I am not sure whether it is best learned from this book. 

\vspace{8pt}

\textbf{Keiji Oguiso. "代数曲線論" Asakura Shoten  (in Japanese)}\vspace{-6pt}

Difficulty: ★★ \vspace{-6pt} 

Comment: If you want to learn algebraic curves and Riemann surfaces from an algebraic-geometric viewpoint, I recommend this book. It is very clear and quite good. You should read it. 
I think this is not a book on algebraic geometry but a book on complex geometry. 
The proofs concerning cohomology are quite analytic. 

Students who do not plan to enter a master's program may read this book in a fourth-year undergraduate seminar. If you plan to enter a master's program, I strongly recommend reading it, even as a supplementary text.
\vspace{8pt}

\textbf{Jean-Pierre Demailly. "Complex Analytic and Differential Geometry"} \vspace{-6pt}

Difficulty: ★★$\sim$★★★★  \vspace{-6pt}

Available online: \url{https://www-fourier.univ-grenoble-alpes.fr/~demailly/manuscripts/agbook.pdf}\vspace{-6pt}

Comment: This book is long. I do not recommend reading it in a seminar, because it is long and contains many rather specialized topics. Chapter 1 is fine, but in Chapter 2 on analytic spaces and Chapter 3 on currents, about half of the material is for specialists. These parts are too difficult. Some portions can be read, but my impression is that specialized material is inserted throughout the book. This is a book that specialists use when studying or citing results.
\vspace{8pt}

\textbf{Phillip Griffiths, Joseph Harris. "Principles of Algebraic Geometry" Wiley Classics Library.}  \vspace{-6pt} 

Difficulty: ★★$\sim$★★★ \vspace{-6pt} 

Comment: This book is long. I do not recommend reading it in a seminar, because it is long. One should not expect to read 800 pages in one year. It should be used like a dictionary.
\vspace{8pt}

\textbf{Hiroshi Konno. "微分幾何学" University of Tokyo Press  (in Japanese)}  \vspace{-6pt} 

Difficulty: ★★$\sim$★★★★ \vspace{-6pt} 

Comment: I think this book covers almost all basic topics in differential geometry. Among books of ordinary length, there are few that contain this much material. It is an excellent book.

However, I do not recommend reading it in a seminar. This is because the arguments and proofs are truly concise and contain no waste. Yet, no matter how many times I read it, for some reason it does not remain in my memory. So it is probably best to read this book after studying from other books.

\subsubsection{For master's seminars}

\textbf{Jean-Pierre Demailly. "Analytic Methods in Algebraic Geometry"  International Press of Boston Inc} \vspace{-6pt}

Difficulty: ★★★$\sim$★★★★ \vspace{-6pt} 

Available online: \url{https://www-fourier.univ-grenoble-alpes.fr/~demailly/manuscripts/analmeth_book.pdf}\vspace{-6pt}

Comment: Master's students in Professor Takayama's research group at the University of Tokyo, including myself, read this book in one year.  This book is not easy to read, and it is difficult. However, after reading it, one can start reading papers from the second year of the master's program.

Professor Takayama's optimized course is ``Chapters 1--6, then Chapters 11--19.'' I also think this is good. Roughly speaking, the contents are as follows.
\vspace{-8pt}
\begin{itemize}[left=0pt]
  \setlength{\parskip}{0cm} % 段落間
  \setlength{\itemsep}{0cm} % 項目間
\item Chapters 1 and 2 consist mainly of definitions, so I think it is enough to read them lightly. One would like to finish these chapters quickly.
\item Chapters 3 and 4 review complex geometry. Here, singular Hermitian metrics, which are central to this field, appear. Note that Professor Demailly uses his own notation for connections and curvature.
\item Chapter 5 treats $L^2$ estimates. If it is difficult to understand, it may be better to consult Professor Demailly's lecture notes on $L^2$ estimates. This chapter introduces multiplier ideal sheaves, which are an important tool.
\item Chapter 6  is the most interesting part, where algebraic-geometric terminology is defined in terms of singular Hermitian metrics.
\item Chapters 11 and 12 treat vector bundles. If you specialize in the $L^2$ theory of several complex variables, you should read this part carefully.
\item Chapter 13 is on the Ohsawa--Takegoshi extension theorem. This is a breakthrough in the field and should be read.
\item Chapter 14 gives applications of the Ohsawa--Takegoshi extension theorem, including Demailly's approximation theorem and the contents of Demailly--Kollar  01. Demailly's approximation is often used, so this chapter should be read.
\item Chapter 15 contains the material of Demailly--Ein--Lazarsfeld 00. This prepares for Chapter 19.
\item Chapter 16 contains the material of Demailly--Peternell--Schneider 01. If you specialize in vanishing theorems, you may read this chapter.
\item Chapter 17 concerns the invariance of plurigenera due to Siu98, Siu03, and P\u{a}un08. It is related to singular Hermitian metrics on vector bundles, so it should be read.
\item Chapter 18 contains the material of Demailly--P\u{a}un 04. When I was a master's student, I thought that it was merely a K\"ahler version of the Nakai--Moishezon criterion. But I now realize how good this theorem is. In the K\"ahler case, there may be no subvarieties at all, so it is not at all obvious that such a theorem should hold.
\item Chapter 19 contains the material of Boucksom--Demailly--P\u{a}un--Peternell 13. It is related to Professor Boucksom's Ph.D thesis. The result is interesting. The statement that $K_X$ is psef if and only if $X$ is nonuniruled is now known, by a result of Professor Wenhao Ou in 2025, to hold also in the K\"ahler case.
\end{itemize}
It may be possible to omit some parts of Chapters 11--19. For example, if you want to do algebraic work, you may omit the latter half of Chapter 14 through Chapter 16; if you want to do analytic work, you may omit Chapters 18 and 19. Personally, I find it unfortunate that Demailly--Peternell--Schneider 94 is not included. That paper would also deserve to be included in a textbook.

My comments on the other chapters are as follows.
\vspace{-8pt}
\begin{itemize}[left=0pt]
  \setlength{\parskip}{0cm} % 段落間
  \setlength{\itemsep}{0cm} % 項目間
\item The parts on the Fujita conjecture in Chapter 7 and Matsusaka's big theorem in Chapter 10 are interesting, but the proofs are technical.  I think it is fine to read them later.
\item Chapter 8 on holomorphic Morse inequalities can probably be read roughly. If you want to read it seriously, you may want to read the book of Ma--Marinescu. It seems that Professor Demailly likes holomorphic Morse inequalities.
\item I do not know whether Chapter 9 on Green--Griffiths--Lang is correct.
\item Chapters 20 and 21 do not need to be read. However, Tsuji's canonical metric in Chapter 20 should have applications, for example in Song--Tian, and Chapter 21 may contain something interesting.
\end{itemize}
\vspace{8pt}

\textbf{Christopher Hacon, Mihnea Popa, Christian Schnell. "Algebraic fiber spaces over abelian varieties: arounda recent theorem by Cao and P\u{a}un"  Amer. Math. Soc., Providence, RI, 2018.}  \vspace{-6pt} 

Difficulty: ★★★$\sim$★★★★ \vspace{-6pt} 

Comment: Students in Professor Takayama's group after me, especially Professors Inayama and Watanabe, probably read this paper after Demailly's book and studied singular Hermitian metrics on vector bundles. I also think this paper is the easiest to understand. If you read it, you should start from the latter half about singular Hermitian metrics. I do not understand the first half, which concerns GV-sheaves and the Fourier--Mukai transform.

Hacon--Popa--Schnell 18 is an exposition of Cao--P\u{a}un 17 and P\u{a}un--Takayama 18. I also recommend reading Professor P\u{a}un's survey (\url{https://arxiv.org/abs/1606.00174}).
\vspace{8pt}

\textbf{Gabor Szekelyhidi. "An Introduction to Extremal Kahler Metrics" Graduate Studies in Mathematics, 152. American Mathematical Society}  \vspace{-6pt} 

Difficulty: ★★★$\sim$★★★★ \vspace{-6pt} 

Comment: If you want to learn the Aubin--Yau theorem and Yau's theorem quickly, this is probably the book. The latter half treats extremal metrics and $K$-stability. The definition of test configuration may be slightly old and may look different from the current refined definition. (It is probably the same.) Still, it may be worthwhile to read such a classical definitions, as well as the definition of the Futaki invariant in Chapter 4.

Topics related to $K$-stability have developed rapidly, so this book may now contain older material. Therefore, if you read it in a seminar, you may read only the first half and use another book for the latter half.
\vspace{8pt}

\textbf{Hiraku Nakajima. "非線形問題と複素幾何学" Iwanami Shoten  (in Japanese)}  \vspace{-6pt} 

Difficulty: ★★★$\sim$★★★★ \vspace{-6pt} 

Comment: This is a good book. It treats the existence problem for K\"ahler--Einstein metrics up to around 1999, including the Aubin--Yau theorem, Yau's theorem, and Nadel's multiplier ideal sheaves. It is a good book that presents the state of the field at that time in full detail.

However, this field has developed rapidly over the past twenty years, and the contents of this book have now become classical. Therefore, it may be safer to read it as a supplementary text. If you read up to the Aubin--Yau theorem and Yau's theorem, this book may be a good choice.
\vspace{8pt}

\textbf{S. Boucksom, T. Hisamoto, M. Jonsson. "Uniform K-stability, Duistermaat-Heckman measures and singularities of pairs." Ann. Inst. Fourier (Grenoble)}  \vspace{-6pt} 

Difficulty: ★★★$\sim$★★★★ \vspace{-6pt} 

Available online: \url{http://sebastien.boucksom.perso.math.cnrs.fr/publis.html}

Comment: If one wants to study $K$-stability from a differential-geometric viewpoint, should one first read this paper? I do not know, so please tell me if there are other good books, papers, or surveys. I remember that this paper itself was fairly easy to read. 

\textbf{Jian Song, Ben Weinkove. "Lecture notes on the Kahler-Ricci flow"}  \vspace{-6pt} 

Difficulty: ★★★$\sim$★★★★ \vspace{-6pt} 

Available online: \url{https://arxiv.org/abs/1212.3653}\vspace{-6pt}

Comment: If you want to study the K\"ahler--Ricci flow, is this the book to read? I read it once, but I have forgotten almost all of its contents. When I was a master's and doctoral student, the K\"ahler--Ricci flow was very active, but recently I do not hear about it as much. What happened to the MMP using the K\"ahler--Ricci flow that Tian talked about? If anyone knows, please tell me.

I noticed later that the contents seem to correspond to Chapter 3 of the book \emph{An Introduction to the Kähler-Ricci Flow}, which is introduced below.
\vspace{8pt}

\textbf{"The Ricci flow: techniques and applications. Part I: Geometric aspects. Chapter 2 Discussion of the Kahler-Ricci flow" }  \vspace{-6pt} 

Difficulty: ★★★$\sim$★★★★ \vspace{-6pt} 

Comment: My collaborator Shiyu Zhang told me about this book. It may indeed be a good way to look quickly at the K\"ahler--Ricci flow. 
\vspace{8pt}

\textbf{Shoshichi Kobayashi. "Differential Geometry of Complex Vector Bundles"  Mathematical Society of Japan}  \vspace{-6pt} 

Difficulty: ★★$\sim$★★★★ \vspace{-6pt} 

Available online: \url{https://www.mathsoc.jp/publications/pubmsj/}\vspace{-6pt}

Comment: The first three chapters can be read even by fourth-year undergraduates. 
However, the contents after Chapter 4 are very hard, so I think these contents after Chapter 4 are  for specialists. 
Chapters 4 and 5 are entirely about the Kobayashi--Hitchin correspondence. I benefited from this material in my research. From Chapter 6 onward, the book turns to moduli, and I did not understand it at all. The book itself is so good that I still cite it.
\vspace{8pt}

\textbf{Hajime Tsuji. "複素多様体論講義" SGC Library 94, Science-sha  (in Japanese)}  \vspace{-6pt} 

Difficulty: ★★★$\sim$★★★★ \vspace{-6pt} 

Comment: This book covers material on ``algebraic geometry using several complex variables.'' It is excellent. As the title suggests, it treats broad foundations. This book should be read as a survey rather than as a textbook.
\vspace{8pt}

\textbf{Sebastien Boucksom, Philippe Eyssidieux, Vincent Guedj. "An Introduction to the Kähler-Ricci Flow" Springer}  \vspace{-6pt} 

Difficulty: ★★★$\sim$★★★★ \vspace{-6pt} 

Comment: I have not read this book. I only referred to Chapter 4 in my research on Bott--Chern classes on singular varieties. It seems to treat the K\"ahler--Ricci flow on singular varieties.

\textbf{Vincent Guedj, Ahmed Zeriahi “Degenerate Complex Monge-Ampere Equations”  Eur. Math. Soc.}
 \vspace{-6pt} 

Difficulty: ★★★$\sim$★★★★ \vspace{-6pt} 

Comment: I have not read this book either. At a glance, the analysis seems difficult. See the following reviews of this book.
\vspace{-8pt}
\begin{itemize}[left=0pt]
  \setlength{\parskip}{0cm} % 段落間
  \setlength{\itemsep}{0cm} % 項目間
\item Review by Genki Hosono (in Japanese) \url{https://www.jstage.jst.go.jp/article/sugaku/74/2/74_0742204/_article/-char/ja}
\item Review by Professor Kolodziej \url{https://link.springer.com/article/10.1365/s13291-018-0182-0}
\end{itemize}
\subsection{Textbooks in several complex variables}

If you go into $L^2$ theory, I think you cannot avoid complex geometry, especially vector bundles and singular Hermitian metrics. Therefore, even if you want to study or specialize in $L^2$ theory, I recommend reading books on complex geometry such as Wells's or Demailly's textbook in a fourth-year undergraduate or master's seminar. Nakano's book is also fine.

For other topics in several complex variables, my suggestions are as follows.
\vspace{-8pt}
\begin{tcolorbox}[mybox]
\begin{itemize}[left=0pt]
  \setlength{\parskip}{0cm} % 段落間
  \setlength{\itemsep}{5pt} % 項目間
\item If you want to study Oka manifolds arising from Oka theory, read Forstneric's book.
\item If you want to study or do research on domains in $\C^n$, such as domains of holomorphy, holomorphic convexity, and pseudoconvexity, I think it is better to study with $L^2$ theory or Oka manifolds as a central theme.
\end{itemize}
\end{tcolorbox}

\subsubsection{For fourth-year undergraduates}

\textbf{Shigeo Nakano. "多変数函数論 - 微分幾何学的アプローチ" Asakura Shoten  (in Japanese)}  　\vspace{-6pt} 

Difficulty: ★★ $\sim$★★★ 　\vspace{-6pt} 

Comment: Before I became a first-year master's student, Professor Takayama told me to read a book on several complex variables, and I hurriedly read this book. I recommend it because one can learn complex geometry and the $L^2$ theory of several complex variables at the same time. If you go into $L^2$ theory, it is also fine to read this book in a fourth-year undergraduate seminar. Unfortunately, it is now hard to obtain. It is a good book.

\vspace{8pt}

\textbf{Franc Forstneric. "Stein Manifolds and Holomorphic Mappings
The Homotopy Principle in Complex Analysis" Springer}   \vspace{-6pt} 

Difficulty: ★★ $\sim$★★★★ \vspace{-6pt} 

Comment: This is a book on Oka manifolds. Personally, I think Oka manifolds may be regarded as a new area of several complex variables. 
Since Japan has Professor Kusakabe, who is a top-class researcher on Oka manifolds, I think this field may become important in Japan.
I strongly encourage you to read it. It is fine to read it either in a fourth-year undergraduate seminar or in a master's seminar. 
However, since I have the impression that Oka manifolds are closer to geometry than to analysis, so complex geometry is ultimately necessary.

When I asked Professor Kusakabe, he said that Chapters 1 and 2 contain complex geometry without proofs, and that one should be able to read Chapters 3--7. Therefore, a standard route may be to study complex geometry first and then read Chapters 3--7.
\vspace{8pt}

\textbf{Lars Hörmander. "Introduction to Complex Analysis in Several Variables" North-Holland Mathematical Library}   \vspace{-6pt} 

Difficulty: ★★ $\sim$★★★ \vspace{-6pt} 

Comment: This is a classical masterpiece in several complex variables. My impression is that fourth-year undergraduates who like hard analysis, such as CR manifolds in the style of Professor Hirachi's research group at the University of Tokyo, read this book.
\vspace{8pt}

\textbf{Takeo Ohsawa. "多変数複素解析" Iwanami Shoten  (in Japanese)}  　\vspace{-6pt} 

Difficulty: ★★ $\sim$★★★★ 　\vspace{-6pt} 

Comment: I have only read a little of this book. If you want to learn domains in $\C^n$, such as domains of holomorphy, holomorphic convexity, and pseudoconvexity, using $L^2$ theory, this may be the Japanese book to read. If you plan to study $L^2$ theory, I think it is better to learn complex geometry as your main topic and read this book as a supplementary text. An expanded edition has also appeared, adding material on the optimal estimates of Ohsawa--Takegoshi due to Guan--Zhou and Blocki. Professor Ohsawa's books contain humor throughout and are enjoyable to read.
\vspace{8pt}

\textbf{Takeo Ohsawa. "複素解析幾何とディーバー方程式 " Baifukan  (in Japanese)}  　\vspace{-6pt} 

Difficulty: ★★ $\sim$★★★★ 　\vspace{-6pt} 

Comment: I have also only read a little of this book. It is an interesting book, but its contents are broad, so it is difficult for beginners. It is probably better to read it after learning the basics. When I read it now, I am constantly surprised that it contains so much material.
\vspace{8pt}

\textbf{Takushiro Ochiai, Junjiro Noguchi. "幾何学的関数論" Iwanami Shoten  (in Japanese)}  　\vspace{-6pt} 

\textbf{Takushiro Ochiai, Junjiro Noguchi. "Geometric Function Theory in Several Complex Variables" American Mathematical Society}  　\vspace{-6pt} 

Difficulty: ★★ $\sim$★★★ 　\vspace{-6pt} 

Comment: I read only the part on currents in Chapter 3. This is one of the few books that gives a detailed account of currents, so Chapter 3 alone is worth reading. Judging from the table of contents, the book concerns hyperbolicity and Nevanlinna theory. I am personally curious about Nevanlinna theory, so I would like someone to read it.
\vspace{8pt}

\textbf{Kunihiko Kodaira. "Nevanlinna Theory" Springer}  　\vspace{-6pt} 

Difficulty: ★★ 　\vspace{-6pt} 

Comment: This book is based on Professor Kodaira's lectures on Nevanlinna theory and was translated into English by Professor Ohsawa. I am also curious about this book and would like someone to read it. I think it can be read if one knows manifolds and complex analysis.
\vspace{8pt}

\textbf{Junjiro Noguchi. "多変数解析関数論 -学部生へおくる岡の連接定理-" Asakura Shoten (in Japanese)}  　\vspace{-6pt} 

\textbf{Junjiro Noguchi. "岡理論新入門 -多変数関数論の基礎- "  Shokabo (in Japanese)}  　\vspace{-6pt} 

Difficulty: ★★ $\sim$★★★★ 　\vspace{-6pt} 

Comment: I think these books treat several complex variables concerning domains in $\C^n$, especially domains of holomorphy, holomorphic convexity, and pseudoconvexity. I do not recommend reading them in a seminar, because the methods are too classical. It is better to read them as supplementary texts.

Such a classical theory, namely the theory of domains in $\C^n$ (domains of holomorphy, holomorphic convexity, and pseudoconvexity), has been reinterpreted through $L^2$ theory and Oka manifolds. And $L^2$ theory and Oka manifolds are easier to study as research topics. Therefore, for students who want to become researchers, I  recommend the latter.

However, I should emphasize that the classical theory itself is very good. In particular, if you treat singular varieties, the classical theory may  be better.
 \vspace{8pt}

\textbf{Hans Grauert, Reinhold Remmert. "Theory of Stein Spaces" Springer }  　\vspace{-6pt} 

Difficulty: ★★ $\sim$★★★★ 　\vspace{-6pt} 

Comment: This is a book on Stein spaces. The contents themselves are very interesting, and reading it now, I am constantly surprised that one can prove so much for analytic spaces without assuming normality, reducedness, or irreducibility. I also read it from time to time. However, the methods are too classical, so I do not recommend reading it in a seminar. It is better to read it as a supplementary text.
\vspace{8pt}

\textbf{Shin Hitotsumatsu. "多変数解析函数論" Baifukan (in Japanese)}  　\vspace{-6pt} 

Difficulty: ★★ $\sim$★★★★ 　\vspace{-6pt} 

Comment: This book also treats very classical several complex variables. Since the methods are too classical, I do not recommend reading it in a seminar. It is better to read it as a supplementary text. However, the final appendix, ``A short history and outlook of several complex variables,'' should be read. This appendix alone gives the book value.
\vspace{8pt}

\textbf{Gerd Fischer. "Complex Analytic Geometry" Springer}  　\vspace{-6pt} 

Difficulty: ★★ $\sim$★★★★ 　\vspace{-6pt} 

Comment: This is a book on analytic spaces. It is extremely useful in research, because it translates algebraic-geometric terminology into complex-analytic terminology. I think this is easier to visualize than scheme theory. It is written carefully and is a good book. However, I do not recommend reading it in a seminar, because it is better read after learning scheme theory.

That said, by reading this book, one can understand complex analytic spaces without understanding scheme theory. In other words, this book is sufficient for doing algebraic geometry over the complex numbers. Therefore, it may be possible to aim for complex algebraic geometry through this book. However, unlike scheme theory, proofs in analytic spaces are quite involved, and one should be prepared for that. At some point, one may feel that schemes are easier.
\vspace{8pt}

\textbf{Satoru Igari. "実解析入門" Iwanami Shoten (in Japanese)}  　\vspace{-6pt} 

Difficulty: ★ $\sim$★★ 　\vspace{-6pt} 

Comment: If you do not understand distributions at all, I would like you to read this book. When I read Nakano's book, I read this book in parallel. I recommend reading it as a supplementary reference book.

\subsubsection{For master's students}

\textbf{Hans Grauert, Thomas Peternell, Reinhold Remmert. "Several Complex Variables VII -Sheaf Theoretical Methods in Complex Analysis-" Springer}  　\vspace{-6pt} 

Difficulty: ★★ $\sim$★★★★ 　\vspace{-6pt} 

Comment: This is an excellent book. If you want to quickly learn several complex variables up to around 1990, from classical several complex variables to $L^2$ theory, this is the book. I am also reading it now. It should be read as a survey rather than as a textbook. Among the many chapters, the parts written by Professor Peternell are especially interesting.
\vspace{8pt}

\textbf{Takeo Ohsawa.  "$L^2$ Approaches in Several Complex Variables
Towards the Oka–Cartan Theory with Precise Bounds" Springer}  　\vspace{-6pt} 

Difficulty: ★★★ $\sim$ ★★★★  　\vspace{-6pt} 

Comment: By reading this book, one can cover $L^2$ theory up to around 2015, especially up to the optimal estimates of the Ohsawa--Takegoshi theorem due to Guan--Zhou and Blocki. This should also be read as a survey rather than as a textbook. See also Professor Matsumura's review in Japanese (\url{https://www.jstage.jst.go.jp/article/sugaku/72/3/72_0723310/_article/-char/ja/}).
\vspace{8pt}

\textbf{Sebastien Boucksom. "Singularities of plurisubharmonic functions and multiplier ideals"}  　\vspace{-6pt} 

Difficulty: ★★★ $\sim$ ★★★★  　\vspace{-6pt} 

Available online: \url{http://sebastien.boucksom.perso.math.cnrs.fr/notes/L2.pdf}\vspace{-6pt}

Comment: These are lecture notes by Professor Boucksom. I strongly recommend reading them, even as a supplementary text. The final part on valuations is truly interesting.
\vspace{8pt}

\textbf{Jean-Pierre Demailly. "$L^2$ estimates for the $\overline{\partial}$-operator on complex manifolds"}  　\vspace{-6pt} 

Difficulty: ★★★ $\sim$ ★★★★  　\vspace{-6pt} 

Available online: \url{https://www-fourier.univ-grenoble-alpes.fr/~demailly/manuscripts/estimations_l2.pdf}\vspace{-6pt}

Comment: These are lecture notes by Professor Demailly. They are useful as a reference when reading the first half, Chapters 1--6, of \emph{Analytic Methods in Algebraic Geometry}.
\vspace{8pt}

\textbf{Bo Berndtsson. "An Introduction to things $\bar\partial$."}  　\vspace{-6pt} 

Difficulty: ★★★ $\sim$ ★★★★  　\vspace{-6pt} 

Available online: \url{https://www.math.chalmers.se/~bob/7nynot.pdf}\vspace{-6pt}

Comment: These are lecture notes by Professor Berndtsson. If Demailly's \emph{Analytic Methods in Algebraic Geometry} is too hard, it may be good to read these notes. 
\vspace{8pt}

\textbf{Shoshichi Kobayashi "Hyperbolic Complex Spaces" Springer}  　\vspace{-6pt} 

Difficulty: ★★★ $\sim$ ★★★★  　\vspace{-6pt} 

Comment: This book contains material on hyperbolicity, such as Kobayashi hyperbolicity and Brody hyperbolicity, from the basics to applications. I looked at it recently and was surprised by the number of techniques involved. For example, in the part on rational maps of hyperbolic manifolds in Chapter 6 (Section 6.6), the argument uses the Douady space and finally applies Miyaoka--Mori to produce a rational curve and derive a contradiction. This is quite impressive. The Kobayashi--Ochiai theorem (the version for varieties of general type) is also contained in this book.
\vspace{8pt}

\subsection{Other books. Books that may be chosen by students graduating after the fourth undergraduate year}
\label{subsec-seminar-23}

The following books may be chosen in a fourth-year undergraduate seminar only by students who will graduate after the fourth year and do not plan to enter a master's program in mathematics.

\subsubsection{Algebra}

\textbf{Michael F. Atiyah, I. G. MacDonald. "Introduction to commutative algebra" CRC Press}  　\vspace{-6pt} 

Difficulty: ★$\sim$★★ 　\vspace{-6pt} 

Comment: If you want to learn the commutative algebra needed for algebraic geometry quickly, this is the book. The exercises are interesting.
If you read it including the exercises, it can take one year.

\vspace{8pt}

\textbf{Yujiro Kawamata. "射影空間の幾何学" Asakura Shoten (in Japanese)}  　\vspace{-6pt} 

Difficulty: ★$\sim$★★ 　\vspace{-6pt} 

Comment: My impression is that this book contains material corresponding to Chapter 1 of Hartshorne. It probably treats classical algebraic geometry of the nineteenth and early twentieth centuries, before Grothendieck's schemes appeared. Famous theorems in plane geometry are written in the language of projective geometry.

The contents are quite interesting. If you do research in algebraic geometry, you cannot avoid abstract arguments involving schemes and sheaves, but if you only want to get a feeling for algebraic geometry, this book is enjoyable enough.
\vspace{8pt}

\textbf{Takeshi Hirai. "線形代数と群の表現I, II" Asakura Shoten (in Japanese)}  　\vspace{-6pt} 

Difficulty: ★$\sim$★★ 　 　\vspace{-6pt} 

Comment: I only glanced at these books, but they explain the basics of group theory through representation theory clearly. If you became lost in algebra in your third undergraduate year, you may study from these books. 

\vspace{8pt}


\textbf{David A Cox, John Little, Donal O'Shea. "Ideals, Varieties, and Algorithms: An Introduction to Computational Algebraic Geometry and Commutative Algebra" Springer}  　\vspace{-6pt} 

Difficulty: ★★$\sim$★★★\vspace{-6pt} 

Comment: When I participated in an internship at NEC, someone asked me, ``your specialty is algebra, right? Do you know about Gr\"obner bases?'' I answered that I did not know them at all. I have been curious about Gr\"obner bases, but I have not studied them so far, so I would like to study them here.
\vspace{8pt}

\textbf{Alexandru Dimca. "Hyperplane Arrangements" Springer}  　\vspace{-6pt} 

\textbf{Peter Orlik , Hiroaki Terao "Arrangements of Hyperplanes" Springer}  　\vspace{-6pt} 

Difficulty: ★★$\sim$★★★\vspace{-6pt} 

Comment: I am somewhat curious about hyperplane arrangements. When I asked Professor Yoshinaga for recommended books, he told me about these two. Dimca is newer, while Orlik--Terao is older but still contains original material.
\vspace{8pt}

\textbf{Reinhard Diestel. "Graph Theory" Springer}  　\vspace{-6pt} 

Difficulty: ★★\vspace{-6pt} 

Comment: Theorems in graph theory are interesting, but I always forget the statements and proofs.
\vspace{8pt}


\textbf{Jean-Pierre Serre. "A Course in Arithmetic" GTM, volume 7. Springer}  　\vspace{-6pt} 

Difficulty: ★★ $\sim$★★★\vspace{-6pt} 

Comment: I am curious about this book and would like someone to read it. I am not good at number theory, so I would like someone to teach me.
\vspace{8pt}


\textbf{James E. Humphreys ."Introduction to Lie Algebras and Representation Theory" Springer} \vspace{-6pt} 

\textbf{Jean-Pierre Serre. "Complex Semisimple Lie Algebras" Springer}  　\vspace{-6pt} 

Difficulty: ★★ $\sim$★★★\vspace{-6pt} 

Comment: Lie algebras are something I have long intended to study someday, but still have not studied. I would like someone to read these books. 

\vspace{8pt}

\subsubsection{Geometry}

\textbf{Shoshichi Kobayashi. "曲線と曲面の微分幾何" Shokabo (in Japanese)}  　\vspace{-6pt} 

\textbf{Reiko Miyaoka. "曲線と曲面の現代幾何学" Iwanami Shoten (in Japanese)}  　\vspace{-6pt} 

\textbf{Hiroyuki Tasaki. "曲線・曲面の微分幾何" Kyoritsu Shuppan (in Japanese)}  　\vspace{-6pt} 

Difficulty: ★$\sim$★★ 　\vspace{-6pt} 

Comment: I have forgotten elementary differential geometry, so I would like to recall it here.

\vspace{8pt}

\textbf{Junichi Inoguchi. "はじめて学ぶリー群 線型代数から始めよう" Gendai Sugakusha (in Japanese)}  　\vspace{-6pt} 

Difficulty: ★$\sim$★★ 　\vspace{-6pt} 

Comment: I have not read this book, but it is apparently readable if one knows linear algebra. \vspace{8pt}

\textbf{Shingo Murakami. "連続群論の基礎 " Asakura Shoten (in Japanese)}  　\vspace{-6pt} 

Difficulty: ★★ 　\vspace{-6pt} 

Comment: After I wrote, as above, ``Please tell me if there is a good textbook on Lie groups,'' someone told me that this book is good. It is indeed easy to understand. Another frequently mentioned textbook on Lie groups is \emph{リー群と表現論} by Toshiyuki Kobayashi and Toshio Oshima in Iwanami Shoten (in Japanese).
\vspace{8pt}

\textbf{J. Matousek. "Using the Borsuk-Ulam Theorem: Lectures on Topological Methods in Combinatorics and Geometry" Springer}  　\vspace{-6pt} 

Difficulty: ★ $\sim$★★\vspace{-6pt} 

Comment: Do you know the ham sandwich theorem? The ham sandwich theorem says that, for a given ham-and-cheese sandwich in real three-dimensional space, there exists a cut that divides the amounts of ham, cheese, and bread each exactly in half. This follows from the Borsuk--Ulam theorem. This book focuses on the Borsuk--Ulam theorem and treats its applications.

\vspace{8pt}

\textbf{J. W. Milnor. "Topology from the differentiable viewpoint" Princeton University Press}  　\vspace{-6pt} 

Difficulty: ★★ \vspace{-6pt} 

Comment: Do you know that there must be a point on the Earth where no wind is blowing? This follows from the Poincare--Hopf theorem. This book treats the relation between differential-geometric objects, such as vector fields, and topological objects, such as the Euler number. Since one cannot read this book without knowing manifolds, it may be good to read a book on manifolds in the first semester and then read this book in the second semester.
\vspace{8pt}

\textbf{Kyoji Kawazumi. "トポロジーの基礎 上・下" University of Tokyo Press (in Japanese)}  　\vspace{-6pt} 

Difficulty: ★ $\sim$★★★\vspace{-6pt} 

Comment: This is probably a book based on the homology course that Professor Kawazumi taught to third-year undergraduates at the University of Tokyo in the second semester. Judging from the table of contents, the first volume contains third-year undergraduate material, and the second volume contains fourth-year undergraduate material. The book has a good reputation (see also Professor Shinichiro Matsuo's page). I think it will become a new standard textbook in topology. Incidentally, I attended that course.

Personally, I would like someone to read it and teach me. The reason is that I have finally reached a point where I can no longer avoid topology in my research. Cup products, excision, and Borel--Moore homology have appeared in my research, and terminology related to topology, such as CW complexes and simplicial sets, also appeared in a seminar on condensed mathematics.
\vspace{8pt}

\textbf{Mikiya Masuda. "代数的トポロジー" Asakura Shoten (in Japanese)}  　\vspace{-6pt} 

Difficulty: ★ $\sim$★★★\vspace{-6pt} 

Comment: This book summarizes the contents of a third-year undergraduate course in algebraic topology, with some additional material. It also contains the Borsuk--Ulam theorem, the Poincare--Hopf theorem, Brouwer's fixed point theorem, and so on. It is very concise. Since it develops the subject in an algebraic style, I personally prefer it. However, people who like geometry may find it too abstract and hard to read.
\vspace{8pt}

\textbf{Colin C. Adams. "The Knot Book: An Elementary Introduction to the Mathematical Theory of Knots" American Mathematical Society}  　\vspace{-6pt} 

Difficulty: ★★ $\sim$★★★\vspace{-6pt} 

Comment: This is not my specialty at all. I looked at it briefly and thought that the figures were beautiful.
\vspace{8pt}

\textbf{Atsushi Kasue. "リーマン幾何学" Baifukan (in Japanese)}  　\vspace{-6pt} 

Difficulty: ★★ $\sim$★★★\vspace{-6pt} 

Comment: Riemannian geometry is another subject that I have long intended to study someday but still have not studied. Therefore, I would like someone to read it.

\vspace{8pt}

\textbf{Barrett O'Neill. "Semi-Riemannian Geometry With Applications to Relativity" Academic Press}  　\vspace{-6pt} 

Difficulty: ★★ $\sim$★★★\vspace{-6pt} 

Comment: When I asked people in geometric analysis at Osaka University, ``Are there any good textbooks on Riemannian geometry?'', they said, ``Well, \emph{Semi-Riemannian Geometry} is good, but...'' They also recommended Professor Kasue's book mentioned above.

\subsubsection{Others}


\textbf{Martin Aigner, G\"unter M. Ziegler. "Proofs from THE BOOK" Springer}  　\vspace{-6pt} 

Difficulty: ★$\sim$★★★ 　\vspace{-6pt} 

Comment: This book collects beautiful proofs in the spirit of Erd\H{o}s. Reading it makes one feel as if one's mind has become sharper. I think it can be read if one knows the material up to the second undergraduate year plus a little more.
\vspace{8pt}

\textbf{Yuichi Ike, E.G. Escolar, Ippei Obayashi, and Shizuo Kaji. "位相的データ解析から構造発見へ パーシステントホモロジーを中心に" Science-sha (in Japanese)}  　\vspace{-6pt} 

Difficulty: ★$\sim$★★★\vspace{-6pt} 

Comment: I am curious about topological data analysis (TDA). I was unsure which book to choose, but I recommend this one because I think that a book written by Professor Ike, who is a mathematical researcher, should be mathematically reliable. I have a rough understanding of the mathematical definitions, but I am curious about the question, ``How are the algorithms constructed?'' I would like someone to read it.

Professor Ike gave an intensive course at Osaka University in 2025. Notes for that course are available on his website (\url{https://sites.google.com/view/yuichi-ike/teaching}), and they may also be useful.
\vspace{8pt}

\textbf{S. Mac Lane "Categories for the Working Mathematician" Springer}  　\vspace{-6pt} 

\textbf{T. Leinster "Basic Category Theory" Cambridge University Press}  　\vspace{-6pt} 

Difficulty: ★$\sim$★★★ \vspace{-6pt} 

Comment: Category theory is often assumed before one realizes it. Students who plan to go into computer science may read it in a fourth-year undergraduate seminar, and afterwards they should become interested in Haskell, which would please one of my classmates.
 If you plan to enter a master's program in mathematics and want to study foundations, topos theory, or other highly abstract mathematics, it is also fine to read category theory in a fourth-year seminar. However, in Japan this field is often included in computer science, so if you enter a master's program, you should also consider going to a computer science department or going abroad.

\vspace{8pt}

\textbf{J.A. Buchmann "Introduction to Cryptography" Springer}  　\vspace{-6pt} 

Difficulty: ★ $\sim$★★\vspace{-6pt} 

Comment: I have wanted to know about cryptography for ten years, but I still have not studied it. I am curious about it.
\vspace{8pt}

\textbf{C.M. Bishop. "Pattern Recognition and Machine Learning" Springer}  　\vspace{-6pt} 

Difficulty: ★ $\sim$★★\vspace{-6pt} 

Comment: I feel that people who want to learn machine learning read this book. However, it has been about twenty years since this book appeared, so I wonder why it is still read. I am curious and would like someone to read it.
\vspace{8pt}

\textbf{Naoyuki Funaki. "確率論" Asakura Shoten (in Japanese)}  　\vspace{-6pt} 

Difficulty: ★★\vspace{-6pt} 

Comment: I once read this book all at once. It was extremely interesting. By now, however, I have forgotten all of its contents. Therefore, I would like to recall them.
\vspace{8pt}

\textbf{David Williams. "Probability with Martingales"  University of Cambridge}  　\vspace{-6pt} 

Difficulty: ★★ $\sim$★★★\vspace{-6pt} 

Comment: When I asked a specialist in probability theory for a recommended textbook, the answer was that this book is easy to read because it summarizes measure theory concisely.

\vspace{8pt}

\textbf{R. Durrett. "Essentials of Stochastic Processes" Springer}  　\vspace{-6pt} 

Difficulty: ★★ $\sim$★★★\vspace{-6pt} 

Comment: I am also curious about probability theory, so I would like someone to read this book.
\vspace{8pt}

\textbf{L. Barreira. "Ergodic Theory, Hyperbolic Dynamics and Dimension Theory." Springer}　\vspace{-6pt} 

\textbf{M. Einsiedler and T. Ward. " Ergodic Theory." Springer }  　\vspace{-6pt} 

Difficulty: ★★ $\sim$★★★★\vspace{-6pt} 

Comment: When I was looking at Professor Kouichi Taira's lecture notes on Lebesgue integration \url{https://sites.google.com/view/the-home-page-of-kouichi-taira/teaching}, I found a problem like the following: for $n \in \mathbb{N}$, let $l_n$ be the leading digit of $2^{n}$. For example, since $2^9=512$, we have $l_9=5$. With what frequency does $l_n=1$ occur? Apparently this can be understood using ergodic theory, and the above books were cited. The answer is said to be in Example 3.4 of Barreira's book. I think I read Barreira's book in a seminar during my master's program, together with Professor Taira and others. In the end, however, I did not really understand it, especially the parts on dynamical systems.
